\SourcePage{603}{606}
\section{Hyperfine structure of molecular levels}

The origin of the hyperfine structure of molecular energy levels is analogous
to that of atomic levels.

In the overwhelming majority of molecules the total electronic spin is zero.
For such molecules, the principal source of hyperfine level splitting is the
quadrupole interaction between the nuclei and the electrons. Naturally, only
nuclei whose spin $i$ differs from 0 and $1/2$ participate in this
interaction, since otherwise the quadrupole moment vanishes.

Because the nuclei move comparatively slowly within a molecule, the operator
of the quadrupole interaction is averaged over a molecular state in two
stages. It must first be averaged over the electronic state with the nuclei
held fixed, and then over the rotation of the molecule.

Consider first a diatomic molecule. After the first averaging stage, the
interaction of each nucleus with the electrons is represented by an operator
proportional to the scalar $\hat Q_{ik}n_in_k$. This scalar is formed from the
operator of the nuclear quadrupole-moment tensor and the unit vector
$\mathbf n$ along the molecular axis, the sole quantity specifying the
orientation of the molecule relative to the direction of the nuclear spin.
Since $\hat Q_{ii}=0$, this operator may be written as
\begin{equation}
 b\hat i_i\hat i_k\left(n_in_k-\frac13\delta_{ik}\right);
 \tag{122.1}
\end{equation}
for a specified value of the projection $i_\zeta$ of the nuclear spin on the
molecular axis, its value is $b[i_\zeta^2-\frac13i(i+1)]$.

On averaging the operator (122.1) over molecular rotation, it is instead
expressed in terms of the operator $\hat{\mathbf K}$ of the conserved
rotational angular momentum. The product $n_in_k$ is averaged by the formula
derived in the problem to \S~29 (with $\mathbf K$ in place of $\mathbf l$),
and the result is
\begin{equation}
 -\frac{b}{(2K-1)(2K+3)}\hat i_i\hat i_k
 \left[\hat K_i\hat K_k+\hat K_k\hat K_i
 -\frac23\delta_{ik}K(K+1)\right].
 \tag{122.2}
\end{equation}
The eigenvalues of this operator are found in the same manner as described
for the operator (121.6).

For a polyatomic molecule, (122.1) is replaced, in general, by an operator of
the form
\begin{equation}
 b_{ik}\hat i_i\hat i_k,
 \tag{122.3}
\end{equation}

\SourcePage{604}{607}
where $b_{ik}$ is a traceless tensor that constitutes a particular
characteristic of the molecule's electronic state. Once it has been averaged
over the molecule's rotation, this tensor is expressed in terms of the total
rotational angular momentum $\mathbf J$ by a formula of the form
\begin{equation}
 \overline{b_{ik}}=b\left[\hat J_i\hat J_k+\hat J_k\hat J_i
 -\frac23J(J+1)\delta_{ik}\right].
 \tag{122.4}
\end{equation}

In principle, the coefficient $b$ can be written in terms of the components of
the tensor $b_{ik}$ referred to the molecule's principal axes of inertia
$\xi$, $\eta$, and $\zeta$. These axes are fixed in the molecule. Consequently,
the components $b_{\xi\xi},\ldots$ characterize the molecule without being
altered by the averaging. To obtain this relation, consider the scalar
$\overline{b_{ik}J_iJ_k}$. Its evaluation using (122.4) gives

\begin{equation}
 \overline{b_{ik}J_iJ_k}
 =bJ(J+1)\left[\frac43J(J+1)-1\right]
 \tag{122.5}
\end{equation}
(the calculation is analogous to the one carried out in the problem to
\S~29). On the other hand, expanding the tensor product along the $\xi$,
$\eta$, $\zeta$ axes yields
\begin{equation}
 \overline{b_{ik}J_iJ_k}
 =b_{\xi\xi}\overline{J_\xi^2}
 +b_{\eta\eta}\overline{J_\eta^2}
 +b_{\zeta\zeta}\overline{J_\zeta^2}.
 \tag{122.6}
\end{equation}
Here it has been used that the mean values of the products
$J_\xi J_\zeta,\ldots$ vanish\footnote{Indeed, in a representation in which
the matrix of one component of $\mathbf J$ (say $J_\zeta$) is diagonal, the
matrices of the products $J_\xi J_\zeta$, $J_\eta J_\zeta$ have nonzero
elements only when the quantum number $k$ changes by 1. The wavefunctions of
the stationary states of an asymmetric top, however, contain functions
$\psi_{Jk}$ whose values of $k$ differ by an even integer (see \S~103).}.
The wave functions of the relevant rotational states of the top allow, in principle, the mean values of $J_\xi^2,\ldots$ to be calculated.
For a symmetric top, in particular, the result has the simple form

\[
 \overline{J_\zeta^2}=k^2,\qquad
 \overline{J_\xi^2}=\overline{J_\eta^2}
 =\frac12[J(J+1)-k^2].
\]

When the nuclear spins are $1/2$, the quadrupole interaction is absent. One
of the principal sources of hyperfine splitting in this case is the direct
magnetic interaction of the nuclear magnetic moments with each other. The
interaction operator of two magnetic moments
$\boldsymbol\mu_1=\mu_1\hat{\mathbf i}_1/i_1$,
$\boldsymbol\mu_2=\mu_2\hat{\mathbf i}_2/i_2$ is
\[
 \frac{\mu_1\mu_2}{i_1i_2r^3}
 [\hat{\mathbf i}_1\hat{\mathbf i}_2
 -3(\hat{\mathbf i}_1\mathbf n)(\hat{\mathbf i}_2\mathbf n)].
\]

\SourcePage{605}{608}
To determine the splitting energy, this operator must be averaged over the
molecular state in the manner described above.

In a molecule containing heavy atoms, the electronic shell also mediates an indirect coupling between the nuclear moments.
Its contribution to the hyperfine splitting can be comparable to that of their direct coupling.
Formally, the indirect contribution arises in second-order perturbation theory from the coupling of a nuclear spin to the electrons.
The results of \S~121 show directly that its magnitude relative to the direct nuclear-moment interaction is of order $(Ze^2/\hbar c)^2$.
For large $Z$, this ratio is of order unity.


Finally, the interaction of a nuclear moment with molecular rotation also
makes a definite contribution to the hyperfine splitting of molecular
levels. As a moving system of charges, a rotating molecule produces a
magnetic field. This field can be calculated from the standard formulas of
electrodynamics for the given current density
$\mathbf j=\rho(\boldsymbol\Omega\times\mathbf r)$, where $\rho$ is the charge
density of the electrons and nuclei in the molecule at rest, and
$\boldsymbol\Omega$ is its angular velocity of rotation. The level splitting
is the energy of the nuclear magnetic moment in this field; the components of
the molecule's angular velocity must here be expressed in terms of the
components of its angular momentum (cf. \S~103).
